% ═══════════════════════════════════════════════════════════════════════════════
\section{Praxisbeispiele}
\label{sec:beispiele}
% ═══════════════════════════════════════════════════════════════════════════════

Dieses Kapitel demonstriert den UQTL-Standard anhand vollständiger,
realistischer Datenbankszenarien. Für jedes Beispiel wird zunächst das
Datenbankschema beschrieben, dann die Abfrage in allen drei Quellsprachen
formuliert, und schließlich das vom Standard definierte normalisierte
SQL-Ergebnis angegeben.

\subsection{Datenmodell für die Beispiele}
\label{subsec:datenmodell}

Alle Beispiele basieren auf folgendem vereinfachten E-Commerce-Schema:

\begin{lstlisting}[style=sqlstyle, caption={Datenbankschema für Beispiele},
  label=lst:schema]
CREATE TABLE customers (
    id          INTEGER PRIMARY KEY,
    name        VARCHAR(200) NOT NULL,
    email       VARCHAR(200) UNIQUE,
    country     VARCHAR(50),
    active      BOOLEAN NOT NULL DEFAULT TRUE,
    created_at  TIMESTAMP NOT NULL,
    tier        VARCHAR(20)  -- 'bronze', 'silver', 'gold', 'platinum'
);

CREATE TABLE orders (
    id          INTEGER PRIMARY KEY,
    customer_id INTEGER REFERENCES customers(id),
    order_date  DATE NOT NULL,
    status      VARCHAR(30),  -- 'pending','confirmed','shipped','delivered'
    total       DECIMAL(12,2) NOT NULL
);

CREATE TABLE order_items (
    id          INTEGER PRIMARY KEY,
    order_id    INTEGER REFERENCES orders(id),
    product_id  INTEGER REFERENCES products(id),
    quantity    INTEGER NOT NULL,
    unit_price  DECIMAL(10,2) NOT NULL
);

CREATE TABLE products (
    id          INTEGER PRIMARY KEY,
    name        VARCHAR(200) NOT NULL,
    category    VARCHAR(100),
    price       DECIMAL(10,2) NOT NULL,
    stock       INTEGER NOT NULL DEFAULT 0
);
\end{lstlisting}

% ─────────────────────────────────────────────────────────────────────────────
\subsection{Beispiel 1: Einfache Filterung und Projektion}
\label{subsec:bsp1}
% ─────────────────────────────────────────────────────────────────────────────

\textbf{Aufgabe:} Alle aktiven Kunden aus Deutschland, nur Name und E-Mail,
alphabetisch sortiert.

\begin{lstlisting}[style=pythonstyle, caption={Beispiel 1 -- Python (SQLAlchemy 2.0)},
  label=lst:bsp1_py]
from sqlalchemy import select
from models import Customer

stmt = (
    select(Customer.name, Customer.email)
    .where(Customer.active == True)
    .where(Customer.country == "DE")
    .order_by(Customer.name.asc())
)
results = session.execute(stmt).all()
\end{lstlisting}

\begin{lstlisting}[style=javastyle, caption={Beispiel 1 -- Java (JPA Criteria API)},
  label=lst:bsp1_java]
CriteriaBuilder cb = em.getCriteriaBuilder();
CriteriaQuery<Object[]> cq = cb.createQuery(Object[].class);
Root<Customer> c = cq.from(Customer.class);

cq.multiselect(c.get("name"), c.get("email"))
  .where(
      cb.and(
          cb.isTrue(c.get("active")),
          cb.equal(c.get("country"), "DE")
      )
  )
  .orderBy(cb.asc(c.get("name")));

List<Object[]> results = em.createQuery(cq).getResultList();
\end{lstlisting}

\begin{lstlisting}[style=csharpstyle, caption={Beispiel 1 -- C\# (EF Core / LINQ)},
  label=lst:bsp1_cs]
// Query Syntax
var results = from c in context.Customers
              where c.Active == true && c.Country == "DE"
              orderby c.Name ascending
              select new { c.Name, c.Email };

// Equivalent Method Syntax
var results = context.Customers
    .Where(c => c.Active == true && c.Country == "DE")
    .OrderBy(c => c.Name)
    .Select(c => new { c.Name, c.Email });
\end{lstlisting}

\begin{successbox}
\textbf{UQTL-Normiertes SQL (identisch für alle drei Quellsprachen):}
\begin{lstlisting}[style=sqlstyle, numbers=none]
SELECT t1.name, t1.email
FROM customers AS t1
WHERE t1.active = TRUE
  AND t1.country = 'DE'
ORDER BY t1.name ASC
\end{lstlisting}
\end{successbox}

% ─────────────────────────────────────────────────────────────────────────────
\subsection{Beispiel 2: Aggregation mit GROUP BY und HAVING}
\label{subsec:bsp2}
% ─────────────────────────────────────────────────────────────────────────────

\textbf{Aufgabe:} Anzahl bestätigter Bestellungen und Gesamtumsatz pro Kunde,
nur Kunden mit mehr als 3 Bestellungen, absteigend nach Umsatz.

\begin{lstlisting}[style=pythonstyle, caption={Beispiel 2 -- Python},
  label=lst:bsp2_py]
from sqlalchemy import select, func

stmt = (
    select(
        Order.customer_id,
        func.count(Order.id).label("order_count"),
        func.sum(Order.total).label("total_revenue")
    )
    .where(Order.status == "confirmed")
    .group_by(Order.customer_id)
    .having(func.count(Order.id) > 3)
    .order_by(func.sum(Order.total).desc())
)
\end{lstlisting}

\begin{lstlisting}[style=javastyle, caption={Beispiel 2 -- Java},
  label=lst:bsp2_java]
CriteriaBuilder cb = em.getCriteriaBuilder();
CriteriaQuery<Object[]> cq = cb.createQuery(Object[].class);
Root<Order> o = cq.from(Order.class);

Expression<Long>    cnt = cb.count(o.get("id"));
Expression<Double>  rev = cb.sum(o.<Double>get("total"));

cq.multiselect(o.get("customerId"), cnt, rev)
  .where(cb.equal(o.get("status"), "confirmed"))
  .groupBy(o.get("customerId"))
  .having(cb.gt(cnt, 3L))
  .orderBy(cb.desc(rev));
\end{lstlisting}

\begin{lstlisting}[style=csharpstyle, caption={Beispiel 2 -- C\#},
  label=lst:bsp2_cs]
var results = context.Orders
    .Where(o => o.Status == "confirmed")
    .GroupBy(o => o.CustomerId)
    .Where(g => g.Count() > 3)
    .OrderByDescending(g => g.Sum(o => o.Total))
    .Select(g => new {
        CustomerId   = g.Key,
        OrderCount   = g.Count(),
        TotalRevenue = g.Sum(o => o.Total)
    });
\end{lstlisting}

\begin{successbox}
\textbf{UQTL-Normiertes SQL:}
\begin{lstlisting}[style=sqlstyle, numbers=none]
SELECT t1.customer_id,
       COUNT(t1.id)    AS order_count,
       SUM(t1.total)   AS total_revenue
FROM orders AS t1
WHERE t1.status = 'confirmed'
GROUP BY t1.customer_id
HAVING COUNT(t1.id) > 3
ORDER BY SUM(t1.total) DESC
\end{lstlisting}
\end{successbox}

% ─────────────────────────────────────────────────────────────────────────────
\subsection{Beispiel 3: INNER JOIN mit mehreren Bedingungen}
\label{subsec:bsp3}
% ─────────────────────────────────────────────────────────────────────────────

\textbf{Aufgabe:} Name und E-Mail der Kunden, die mindestens eine
gelieferte Bestellung im Wert über 500~EUR hatten, im Jahr 2024.

\begin{lstlisting}[style=pythonstyle, caption={Beispiel 3 -- Python (JOIN)},
  label=lst:bsp3_py]
from sqlalchemy import select, extract

stmt = (
    select(Customer.name, Customer.email)
    .join(Order, Customer.id == Order.customer_id)
    .where(Order.status == "delivered")
    .where(Order.total > 500)
    .where(extract("year", Order.order_date) == 2024)
    .distinct()
)
\end{lstlisting}

\begin{lstlisting}[style=javastyle, caption={Beispiel 3 -- Java (JOIN)},
  label=lst:bsp3_java]
CriteriaBuilder cb = em.getCriteriaBuilder();
CriteriaQuery<Object[]> cq = cb.createQuery(Object[].class);
Root<Customer> c = cq.from(Customer.class);
Join<Customer, Order> o = c.join("orders", JoinType.INNER);

cq.multiselect(c.get("name"), c.get("email"))
  .where(cb.and(
      cb.equal(o.get("status"), "delivered"),
      cb.gt(o.<BigDecimal>get("total"), new BigDecimal("500")),
      cb.equal(
          cb.function("YEAR", Integer.class, o.get("orderDate")),
          2024
      )
  ))
  .distinct(true);
\end{lstlisting}

\begin{lstlisting}[style=csharpstyle, caption={Beispiel 3 -- C\# (JOIN)},
  label=lst:bsp3_cs]
var results = (from c in context.Customers
               join o in context.Orders
                 on c.Id equals o.CustomerId
               where o.Status == "delivered"
                  && o.Total > 500m
                  && o.OrderDate.Year == 2024
               select new { c.Name, c.Email })
              .Distinct();
\end{lstlisting}

\begin{successbox}
\textbf{UQTL-Normiertes SQL:}
\begin{lstlisting}[style=sqlstyle, numbers=none]
SELECT DISTINCT t1.name, t1.email
FROM customers AS t1
INNER JOIN orders AS t2
  ON t1.id = t2.customer_id
WHERE t2.status = 'delivered'
  AND t2.total > 500
  AND EXTRACT(YEAR FROM t2.order_date) = 2024
\end{lstlisting}
\end{successbox}

% ─────────────────────────────────────────────────────────────────────────────
\subsection{Beispiel 4: LEFT OUTER JOIN und NULL-Behandlung}
\label{subsec:bsp4}
% ─────────────────────────────────────────────────────────────────────────────

\textbf{Aufgabe:} Alle Kunden mit Anzahl ihrer Bestellungen (auch solche ohne Bestellungen).

\begin{lstlisting}[style=pythonstyle, caption={Beispiel 4 -- Python (LEFT JOIN)},
  label=lst:bsp4_py]
from sqlalchemy import select, func, outerjoin

stmt = (
    select(
        Customer.id,
        Customer.name,
        func.count(Order.id).label("order_count")
    )
    .outerjoin(Order, Customer.id == Order.customer_id)
    .group_by(Customer.id, Customer.name)
    .order_by(Customer.name)
)
\end{lstlisting}

\begin{lstlisting}[style=javastyle, caption={Beispiel 4 -- Java (LEFT JOIN)},
  label=lst:bsp4_java]
CriteriaBuilder cb = em.getCriteriaBuilder();
CriteriaQuery<Object[]> cq = cb.createQuery(Object[].class);
Root<Customer> c = cq.from(Customer.class);
Join<Customer, Order> o = c.join("orders", JoinType.LEFT);

cq.multiselect(c.get("id"), c.get("name"), cb.count(o.get("id")))
  .groupBy(c.get("id"), c.get("name"))
  .orderBy(cb.asc(c.get("name")));
\end{lstlisting}

\begin{lstlisting}[style=csharpstyle, caption={Beispiel 4 -- C\# (LEFT JOIN)},
  label=lst:bsp4_cs]
var results = from c in context.Customers
              join o in context.Orders
                on c.Id equals o.CustomerId into orders
              from o in orders.DefaultIfEmpty()
              group o by new { c.Id, c.Name } into g
              orderby g.Key.Name
              select new {
                  g.Key.Id,
                  g.Key.Name,
                  OrderCount = g.Count(o => o != null)
              };
\end{lstlisting}

\begin{successbox}
\textbf{UQTL-Normiertes SQL:}
\begin{lstlisting}[style=sqlstyle, numbers=none]
SELECT t1.id, t1.name, COUNT(t2.id) AS order_count
FROM customers AS t1
LEFT OUTER JOIN orders AS t2
  ON t1.id = t2.customer_id
GROUP BY t1.id, t1.name
ORDER BY t1.name ASC
\end{lstlisting}
\end{successbox}

% ─────────────────────────────────────────────────────────────────────────────
\subsection{Beispiel 5: Fensterfunktionen (OVER PARTITION BY)}
\label{subsec:bsp5}
% ─────────────────────────────────────────────────────────────────────────────

\textbf{Aufgabe:} Für jede Bestellung: Rang innerhalb des Kunden nach Bestellwert
(teuerste = Rang 1), laufende Summe der Bestellwerte pro Kunde.

\begin{lstlisting}[style=pythonstyle, caption={Beispiel 5 -- Python (Window Functions)},
  label=lst:bsp5_py]
from sqlalchemy import select, func
from sqlalchemy import over

rnk = func.rank().over(
    partition_by=Order.customer_id,
    order_by=Order.total.desc()
).label("rnk")

run_total = func.sum(Order.total).over(
    partition_by=Order.customer_id,
    order_by=Order.order_date.asc(),
    rows=(None, 0)  # UNBOUNDED PRECEDING to CURRENT ROW
).label("running_total")

stmt = select(
    Order.id,
    Order.customer_id,
    Order.total,
    rnk,
    run_total
).order_by(Order.customer_id, rnk)
\end{lstlisting}

\begin{lstlisting}[style=javastyle, caption={Beispiel 5 -- Java (Window via JPQL / Blaze-Persistence)},
  label=lst:bsp5_java]
// Java JPA hat keine native Window Function API.
// Blaze-Persistence oder native Queries werden benoetigt:
String jpql = """
    SELECT o.id, o.customerId, o.total,
           RANK() OVER (PARTITION BY o.customerId ORDER BY o.total DESC),
           SUM(o.total) OVER (
               PARTITION BY o.customerId
               ORDER BY o.orderDate
               ROWS BETWEEN UNBOUNDED PRECEDING AND CURRENT ROW
           )
    FROM Order o
    ORDER BY o.customerId, 4
    """;
List<Object[]> results = em.createQuery(jpql).getResultList();
\end{lstlisting}

\begin{lstlisting}[style=csharpstyle, caption={Beispiel 5 -- C\# (EF Core Window Functions)},
  label=lst:bsp5_cs]
// EF Core 8+ mit Window Function Unterstuetzung:
var results = context.Orders
    .Select(o => new {
        o.Id,
        o.CustomerId,
        o.Total,
        Rank = EF.Functions.Rank()
            .Over()
            .PartitionBy(o.CustomerId)
            .OrderByDescending(o.Total),
        RunningTotal = EF.Functions.Sum(o.Total)
            .Over()
            .PartitionBy(o.CustomerId)
            .OrderBy(o.OrderDate)
            .RowsBetween(
                EF.Functions.RowsOrRangeBoundary.UnboundedPreceding,
                EF.Functions.RowsOrRangeBoundary.CurrentRow)
    })
    .OrderBy(x => x.CustomerId)
    .ThenBy(x => x.Rank);
\end{lstlisting}

\begin{successbox}
\textbf{UQTL-Normiertes SQL:}
\begin{lstlisting}[style=sqlstyle, numbers=none]
SELECT t1.id,
       t1.customer_id,
       t1.total,
       RANK() OVER (
           PARTITION BY t1.customer_id
           ORDER BY t1.total DESC
       ) AS rnk,
       SUM(t1.total) OVER (
           PARTITION BY t1.customer_id
           ORDER BY t1.order_date ASC
           ROWS BETWEEN UNBOUNDED PRECEDING AND CURRENT ROW
       ) AS running_total
FROM orders AS t1
ORDER BY t1.customer_id ASC, rnk ASC
\end{lstlisting}
\end{successbox}

% ─────────────────────────────────────────────────────────────────────────────
\subsection{Beispiel 6: Unterabfragen (Subqueries)}
\label{subsec:bsp6}
% ─────────────────────────────────────────────────────────────────────────────

\textbf{Aufgabe:} Kunden, deren letzter Bestellwert über dem Durchschnitt
aller Bestellwerte liegt.

\begin{lstlisting}[style=pythonstyle, caption={Beispiel 6 -- Python (Subquery)},
  label=lst:bsp6_py]
from sqlalchemy import select, func

# Unterabfrage: Durchschnittlicher Bestellwert
avg_subq = select(func.avg(Order.total)).scalar_subquery()

# Unterabfrage: Letzte Bestellung je Kunde
last_order = (
    select(Order.customer_id, func.max(Order.order_date).label("last_date"))
    .group_by(Order.customer_id)
    .subquery()
)

# Hauptabfrage
stmt = (
    select(Customer.id, Customer.name)
    .join(last_order, Customer.id == last_order.c.customer_id)
    .join(Order, (Order.customer_id == Customer.id) &
                 (Order.order_date == last_order.c.last_date))
    .where(Order.total > avg_subq)
    .order_by(Customer.name)
)
\end{lstlisting}

\begin{lstlisting}[style=javastyle, caption={Beispiel 6 -- Java (Subquery)},
  label=lst:bsp6_java]
CriteriaBuilder cb = em.getCriteriaBuilder();
CriteriaQuery<Object[]> cq = cb.createQuery(Object[].class);
Root<Customer> c = cq.from(Customer.class);
Join<Customer, Order> o = c.join("orders", JoinType.INNER);

// Scalar Subquery: Durchschnitt
Subquery<Double> avgSub = cq.subquery(Double.class);
Root<Order> so = avgSub.from(Order.class);
avgSub.select(cb.avg(so.<Double>get("total")));

cq.multiselect(c.get("id"), c.get("name"))
  .where(cb.and(
      cb.gt(o.<Double>get("total"), avgSub),
      // letzte Bestellung via correlated subquery ...
      cb.equal(o.get("orderDate"),
          cb.greatest(o.<java.time.LocalDate>get("orderDate")))
  ))
  .orderBy(cb.asc(c.get("name")));
\end{lstlisting}

\begin{lstlisting}[style=csharpstyle, caption={Beispiel 6 -- C\# (Subquery via LINQ)},
  label=lst:bsp6_cs]
var avgTotal = context.Orders.Average(o => o.Total);

var lastOrders = context.Orders
    .GroupBy(o => o.CustomerId)
    .Select(g => new {
        CustomerId = g.Key,
        LastDate   = g.Max(o => o.OrderDate)
    });

var results = from c in context.Customers
              join lo in lastOrders on c.Id equals lo.CustomerId
              join o in context.Orders
                on new { c.Id, lo.LastDate }
                equals new { Id = o.CustomerId, LastDate = o.OrderDate }
              where o.Total > avgTotal
              orderby c.Name
              select new { c.Id, c.Name };
\end{lstlisting}

\begin{successbox}
\textbf{UQTL-Normiertes SQL:}
\begin{lstlisting}[style=sqlstyle, numbers=none]
SELECT t1.id, t1.name
FROM customers AS t1
INNER JOIN orders AS t2
  ON t1.id = t2.customer_id
WHERE t2.order_date = (
    SELECT MAX(t3.order_date)
    FROM orders AS t3
    WHERE t3.customer_id = t1.id
)
  AND t2.total > (
    SELECT AVG(t4.total)
    FROM orders AS t4
)
ORDER BY t1.name ASC
\end{lstlisting}
\end{successbox}

% ─────────────────────────────────────────────────────────────────────────────
\subsection{Beispiel 7: Komplexer Multi-Join mit Aggregation}
\label{subsec:bsp7}
% ─────────────────────────────────────────────────────────────────────────────

\textbf{Aufgabe:} Top-5-Produkte nach Umsatz (Menge $\times$ Einzelpreis),
nur Produkte der Kategorie ,,Elektronik'', mit Anzahl unterschiedlicher Käufer.

\begin{lstlisting}[style=pythonstyle, caption={Beispiel 7 -- Python (Multi-Join + Aggregation)},
  label=lst:bsp7_py]
from sqlalchemy import select, func

stmt = (
    select(
        Product.id,
        Product.name,
        func.sum(OrderItem.quantity * OrderItem.unit_price).label("revenue"),
        func.count(Order.customer_id.distinct()).label("unique_buyers")
    )
    .join(OrderItem, Product.id == OrderItem.product_id)
    .join(Order, OrderItem.order_id == Order.id)
    .where(Product.category == "Elektronik")
    .group_by(Product.id, Product.name)
    .order_by(func.sum(
        OrderItem.quantity * OrderItem.unit_price).desc())
    .limit(5)
)
\end{lstlisting}

\begin{lstlisting}[style=javastyle, caption={Beispiel 7 -- Java (Multi-Join + Aggregation)},
  label=lst:bsp7_java]
CriteriaBuilder cb = em.getCriteriaBuilder();
CriteriaQuery<Object[]> cq = cb.createQuery(Object[].class);
Root<Product> p = cq.from(Product.class);
Join<Product, OrderItem> oi = p.join("orderItems", JoinType.INNER);
Join<OrderItem, Order>   o  = oi.join("order",      JoinType.INNER);

Expression<BigDecimal> revenue =
    cb.sum(cb.prod(oi.<Integer>get("quantity"),
                   oi.<BigDecimal>get("unitPrice")));
Expression<Long> buyers =
    cb.countDistinct(o.get("customerId"));

cq.multiselect(p.get("id"), p.get("name"), revenue, buyers)
  .where(cb.equal(p.get("category"), "Elektronik"))
  .groupBy(p.get("id"), p.get("name"))
  .orderBy(cb.desc(revenue));

List<Object[]> result = em.createQuery(cq)
    .setMaxResults(5)
    .getResultList();
\end{lstlisting}

\begin{lstlisting}[style=csharpstyle, caption={Beispiel 7 -- C\# (Multi-Join + Aggregation)},
  label=lst:bsp7_cs]
var results = context.Products
    .Where(p => p.Category == "Elektronik")
    .Join(context.OrderItems,
          p  => p.Id,
          oi => oi.ProductId,
          (p, oi) => new { p, oi })
    .Join(context.Orders,
          x  => x.oi.OrderId,
          o  => o.Id,
          (x, o) => new { x.p, x.oi, o })
    .GroupBy(x => new { x.p.Id, x.p.Name })
    .Select(g => new {
        g.Key.Id,
        g.Key.Name,
        Revenue      = g.Sum(x => x.oi.Quantity * x.oi.UnitPrice),
        UniqueBuyers = g.Select(x => x.o.CustomerId).Distinct().Count()
    })
    .OrderByDescending(x => x.Revenue)
    .Take(5);
\end{lstlisting}

\begin{successbox}
\textbf{UQTL-Normiertes SQL:}
\begin{lstlisting}[style=sqlstyle, numbers=none]
SELECT t1.id,
       t1.name,
       SUM(t2.quantity * t2.unit_price)   AS revenue,
       COUNT(DISTINCT t3.customer_id)     AS unique_buyers
FROM products AS t1
INNER JOIN order_items AS t2
  ON t1.id = t2.product_id
INNER JOIN orders AS t3
  ON t2.order_id = t3.id
WHERE t1.category = 'Elektronik'
GROUP BY t1.id, t1.name
ORDER BY SUM(t2.quantity * t2.unit_price) DESC
LIMIT 5
\end{lstlisting}
\end{successbox}

% ─────────────────────────────────────────────────────────────────────────────
\subsection{Beispiel 8: EXISTS-Unterabfrage}
\label{subsec:bsp8}
% ─────────────────────────────────────────────────────────────────────────────

\textbf{Aufgabe:} Alle Kunden, die mindestens eine Bestellung mit dem Status
,,shipped'' aber noch keine mit dem Status ,,delivered'' haben.

\begin{lstlisting}[style=pythonstyle, caption={Beispiel 8 -- Python (EXISTS)},
  label=lst:bsp8_py]
from sqlalchemy import select, exists

has_shipped = exists(
    select(Order.id)
    .where(Order.customer_id == Customer.id)
    .where(Order.status == "shipped")
)
has_delivered = exists(
    select(Order.id)
    .where(Order.customer_id == Customer.id)
    .where(Order.status == "delivered")
)

stmt = (
    select(Customer)
    .where(has_shipped)
    .where(~has_delivered)
    .order_by(Customer.name)
)
\end{lstlisting}

\begin{lstlisting}[style=javastyle, caption={Beispiel 8 -- Java (EXISTS)},
  label=lst:bsp8_java]
CriteriaBuilder cb = em.getCriteriaBuilder();
CriteriaQuery<Customer> cq = cb.createQuery(Customer.class);
Root<Customer> c = cq.from(Customer.class);

Subquery<Integer> shipped   = cq.subquery(Integer.class);
Root<Order> os = shipped.from(Order.class);
shipped.select(cb.literal(1))
    .where(cb.and(
        cb.equal(os.get("customerId"), c.get("id")),
        cb.equal(os.get("status"), "shipped")
    ));

Subquery<Integer> delivered = cq.subquery(Integer.class);
Root<Order> od = delivered.from(Order.class);
delivered.select(cb.literal(1))
    .where(cb.and(
        cb.equal(od.get("customerId"), c.get("id")),
        cb.equal(od.get("status"), "delivered")
    ));

cq.select(c)
  .where(cb.and(
      cb.exists(shipped),
      cb.not(cb.exists(delivered))
  ))
  .orderBy(cb.asc(c.get("name")));
\end{lstlisting}

\begin{lstlisting}[style=csharpstyle, caption={Beispiel 8 -- C\# (EXISTS via Any)},
  label=lst:bsp8_cs]
var results = context.Customers
    .Where(c =>
        context.Orders.Any(o => o.CustomerId == c.Id
                             && o.Status == "shipped") &&
       !context.Orders.Any(o => o.CustomerId == c.Id
                             && o.Status == "delivered"))
    .OrderBy(c => c.Name);
\end{lstlisting}

\begin{successbox}
\textbf{UQTL-Normiertes SQL:}
\begin{lstlisting}[style=sqlstyle, numbers=none]
SELECT t1.id, t1.name, t1.email, t1.country,
       t1.active, t1.created_at, t1.tier
FROM customers AS t1
WHERE EXISTS (
    SELECT 1
    FROM orders AS t2
    WHERE t2.customer_id = t1.id
      AND t2.status = 'shipped'
)
  AND NOT EXISTS (
    SELECT 1
    FROM orders AS t3
    WHERE t3.customer_id = t1.id
      AND t3.status = 'delivered'
)
ORDER BY t1.name ASC
\end{lstlisting}
\end{successbox}

% ─────────────────────────────────────────────────────────────────────────────
\subsection{Beispiel 9: Common Table Expressions (WITH)}
\label{subsec:bsp9}
% ─────────────────────────────────────────────────────────────────────────────

\textbf{Aufgabe:} Kunden mit überdurchschnittlichem Umsatz (Top-Tier-Analyse),
unter Verwendung von CTEs für Lesbarkeit.

\begin{lstlisting}[style=pythonstyle, caption={Beispiel 9 -- Python (CTE)},
  label=lst:bsp9_py]
from sqlalchemy import select, func

# CTE 1: Gesamtumsatz je Kunde
customer_revenue = (
    select(
        Order.customer_id,
        func.sum(Order.total).label("total_revenue")
    )
    .group_by(Order.customer_id)
    .cte("customer_revenue")
)

# CTE 2: Durchschnittsumsatz
avg_revenue = (
    select(func.avg(customer_revenue.c.total_revenue).label("avg_rev"))
    .cte("avg_revenue")
)

# Hauptabfrage
stmt = (
    select(Customer.name, Customer.email,
           customer_revenue.c.total_revenue)
    .join(customer_revenue,
          Customer.id == customer_revenue.c.customer_id)
    .where(customer_revenue.c.total_revenue >
           select(avg_revenue.c.avg_rev).scalar_subquery())
    .order_by(customer_revenue.c.total_revenue.desc())
)
\end{lstlisting}

\begin{lstlisting}[style=csharpstyle, caption={Beispiel 9 -- C\# (CTE-Aequivalent)},
  label=lst:bsp9_cs]
// EF Core: CTEs werden als verkettete LINQ-Ausdruecke modelliert
var customerRevenue = context.Orders
    .GroupBy(o => o.CustomerId)
    .Select(g => new {
        CustomerId    = g.Key,
        TotalRevenue  = g.Sum(o => o.Total)
    });

var avgRevenue = customerRevenue.Average(x => x.TotalRevenue);

var results = context.Customers
    .Join(customerRevenue,
          c  => c.Id,
          cr => cr.CustomerId,
          (c, cr) => new { c, cr })
    .Where(x => x.cr.TotalRevenue > avgRevenue)
    .OrderByDescending(x => x.cr.TotalRevenue)
    .Select(x => new {
        x.c.Name,
        x.c.Email,
        x.cr.TotalRevenue
    });
\end{lstlisting}

\begin{successbox}
\textbf{UQTL-Normiertes SQL (mit CTE):}
\begin{lstlisting}[style=sqlstyle, numbers=none]
WITH customer_revenue AS (
    SELECT t2.customer_id,
           SUM(t2.total) AS total_revenue
    FROM orders AS t2
    GROUP BY t2.customer_id
),
avg_revenue AS (
    SELECT AVG(cr.total_revenue) AS avg_rev
    FROM customer_revenue AS cr
)
SELECT t1.name, t1.email, cr.total_revenue
FROM customers AS t1
INNER JOIN customer_revenue AS cr
  ON t1.id = cr.customer_id
WHERE cr.total_revenue > (SELECT avg_rev FROM avg_revenue)
ORDER BY cr.total_revenue DESC
\end{lstlisting}
\end{successbox}

% ─────────────────────────────────────────────────────────────────────────────
\subsection{Beispiel 10: Vollständige Geschäftslogik-Abfrage}
\label{subsec:bsp10}
% ─────────────────────────────────────────────────────────────────────────────

\textbf{Aufgabe:} Monatlicher Umsatzbericht für das Jahr 2024: Monat, Gesamtumsatz,
Anzahl Bestellungen, Anzahl eindeutiger Kunden, Vergleich zum Vormonat
(via LAG-Fensterfunktion).

\begin{lstlisting}[style=pythonstyle, caption={Beispiel 10 -- Python (Business Report)},
  label=lst:bsp10_py]
from sqlalchemy import select, func, extract

monthly = (
    select(
        extract("year",  Order.order_date).label("year"),
        extract("month", Order.order_date).label("month"),
        func.sum(Order.total).label("revenue"),
        func.count(Order.id).label("order_count"),
        func.count(Order.customer_id.distinct()).label("unique_customers")
    )
    .where(extract("year", Order.order_date) == 2024)
    .where(Order.status != "pending")
    .group_by(
        extract("year",  Order.order_date),
        extract("month", Order.order_date)
    )
    .order_by(
        extract("year",  Order.order_date),
        extract("month", Order.order_date)
    )
    .subquery()
)

prev_rev = func.lag(monthly.c.revenue, 1).over(
    order_by=[monthly.c.year, monthly.c.month]
).label("prev_month_revenue")

stmt = select(monthly, prev_rev)
\end{lstlisting}

\begin{lstlisting}[style=csharpstyle, caption={Beispiel 10 -- C\# (Business Report)},
  label=lst:bsp10_cs]
var monthly = context.Orders
    .Where(o => o.OrderDate.Year == 2024
             && o.Status != "pending")
    .GroupBy(o => new {
        Year  = o.OrderDate.Year,
        Month = o.OrderDate.Month
    })
    .Select(g => new {
        g.Key.Year,
        g.Key.Month,
        Revenue         = g.Sum(o => o.Total),
        OrderCount      = g.Count(),
        UniqueCustomers = g.Select(o => o.CustomerId).Distinct().Count()
    })
    .OrderBy(x => x.Year)
    .ThenBy(x => x.Month);

// LAG-Fenster: EF Core 8+
var withLag = monthly
    .Select(m => new {
        m.Year, m.Month, m.Revenue, m.OrderCount, m.UniqueCustomers,
        PrevMonthRevenue = EF.Functions.Lag(m.Revenue, 1)
            .Over()
            .OrderBy(m.Year).ThenBy(m.Month)
    });
\end{lstlisting}

\begin{successbox}
\textbf{UQTL-Normiertes SQL:}
\begin{lstlisting}[style=sqlstyle, numbers=none]
SELECT sub.year,
       sub.month,
       sub.revenue,
       sub.order_count,
       sub.unique_customers,
       LAG(sub.revenue, 1) OVER (
           ORDER BY sub.year ASC, sub.month ASC
       ) AS prev_month_revenue
FROM (
    SELECT EXTRACT(YEAR  FROM t1.order_date) AS year,
           EXTRACT(MONTH FROM t1.order_date) AS month,
           SUM(t1.total)                      AS revenue,
           COUNT(t1.id)                       AS order_count,
           COUNT(DISTINCT t1.customer_id)     AS unique_customers
    FROM orders AS t1
    WHERE EXTRACT(YEAR FROM t1.order_date) = 2024
      AND t1.status != 'pending'
    GROUP BY EXTRACT(YEAR  FROM t1.order_date),
             EXTRACT(MONTH FROM t1.order_date)
    ORDER BY EXTRACT(YEAR  FROM t1.order_date) ASC,
             EXTRACT(MONTH FROM t1.order_date) ASC
) AS sub
\end{lstlisting}
\end{successbox}
